\documentclass[11pt,a4paper]{article}
\usepackage[utf8]{inputenc}
\usepackage[brazil]{babel}
\usepackage{graphicx}
\usepackage{amssymb, amsfonts, amsmath}
\usepackage{color}
\usepackage{float}
\usepackage{enumerate}
%\usepackage{subfigure}
\usepackage[top=2.5cm, bottom=2.5cm, left=1.25cm, right=1.25cm]{geometry}

\newcommand{\sen}{\,\textrm{sen}\,}
\newcommand{\tg}{\,\textrm{tg}\,}

\begin{document}
\pagestyle{empty}

\begin{center}
\begin{tabular}{ccc}
\begin{tabular}{c}
\includegraphics[scale=0.25]{../../biblioteca/imagem/brasao-de-armas-brasil} \\
\end{tabular} & 
\begin{tabular}{c}
Ministério da Educação \\
Universidade Federal dos Vales do Jequitinhonha e Mucuri \\
Faculdade de Ciências Sociais, Aplicadas e Exatas - FACSAE \\
Departamento de Ciências Exatas - DCEX \\
Disciplina: Cálculo Numérico\\
Prof.: Luiz C. M. de Aquino\\
\end{tabular} &
\begin{tabular}{c}
\includegraphics[scale=0.25]{../../biblioteca/imagem/logo-ufvjm} \\
\end{tabular}
\end{tabular}

Aluno(a): \rule{0.5\textwidth}{0.01cm} \quad Data: \rule{0.6cm}{0.01cm} / \rule{0.6cm}{0.01cm} / \rule{1.25cm}{0.01cm}
\end{center}

\begin{center}
 \textbf{Avaliação III}
\end{center}

\textbf{Instruções}
\begin{itemize}
 \item Todas as justificativas necessárias na solução de cada questão devem estar presentes nesta avaliação;
 \item As respostas finais de cada questão devem estar escritas de caneta;
 \item Esta avaliação tem um total de 25,0 pontos.
\end{itemize}

\begin{enumerate}
 
 \item \textbf{[5,0 pontos]} Seja uma função $f$ da qual são conhecidos os valores descritos na tabela abaixo.
\begin{center}
 \begin{tabular}{c|c|c|c}
  $x_i$ & 1 & 1,5 & 2\\ \hline
  $f(x_i)$ & 2 & 3,125 & 6
 \end{tabular}
\end{center}

  Determine o polinômio $p$ que interpola $f$ utilizando duas maneiras:
  \begin{enumerate}
   \item escrevendo $p$ na Forma de Lagrange;
   \item escrevendo $p$ na Forma de Newton.
  \end{enumerate}

  \item \textbf{[6,0 pontos]} Arme os sistemas de equações que permitem determinar uma \textit{spline} natural cúbica que interpole os pontos da tabela abaixo:

   \begin{center}
   \begin{tabular}{c|c|c|c|c}
      $x_i$ & $-1,2$ & $-0,1$ & $0,5$ & $1,4$ \\ \hline
      $y_i$ & $-1$ & $0,2$ & $0,75$ & $-0,15$
   \end{tabular}
   \end{center}

  \item \textbf{[4,0 pontos]} Considere os polinômios:
$$p_0(x) = 1;\, p_1(x) = x;\,p_2(x) = \frac{1}{2}(3x^2 - 1).$$


   \begin{enumerate}
    \item Verifique que $\displaystyle\int_{-1}^1 p_i(x)p_j(x)\,dx = 0$, sempre que $i\neq j$.
    \item Utilize o Método dos Mínimos Quadrados para determinar $\phi(x) = a_0p_0(x) + a_1p_1(x) + a_2p_2(x)$ que melhor se ajusta a função definida por 
          $f(x) = \left(x - \dfrac{1}{2}\right)^4$ no intervalo $[-1,\, 1]$.
   \end{enumerate}


 \item \textbf{[5,0 pontos]} Sobre certa função $f$ são conhecidos os pontos $(x_k,\,f(x_k))$, com $k=0$, $1$, $2$, \ldots, $n$. Suponha que seja 
aplicado o Método dos Mínimos Quadrados para determinar a função $\phi(x) = ag_1(x) + bg_2(x)$ que melhor se ajusta a $f$. Deduza que os coeficientes $a$ e 
$b$ são a solução do sistema de equações:
$$\begin{cases}
   c_{11}a + c_{12}b = d_1 \\
   c_{21}a + c_{22}b = d_2
  \end{cases},
$$
onde $\displaystyle c_{ij} = \sum_{k=0}^n g_i(x_k)g_j(x_k)$ e $\displaystyle d_{i} = \sum_{k=0}^n g_i(x_k)f(x_k)$.

 \item \textbf{[5,0 pontos]} Seja uma função $f$ da qual são conhecidos os pontos $(x_0,\,f(x_0))$ e $(x_1,\,f(x_1))$. Considere que 
$L(x)$ seja o polinômio na Forma de Lagrange que interpola $f$. Além disso, considere que $N(x)$ seja o polinômio na Forma de Newton que 
interpola $f$. Prove que $L(x)$ e $N(x)$ representam um mesmo polinômio. 

\end{enumerate}
\end{document}
